%!TEX root = ../main.tex
\newpage
\chapter{Conclusion}\label{chap:conclusion}

本研究では、物理的に厳密な二媒質屈折モデルを3D Gaussian Splatting (3DGS) の最適化プロセスに直接統合した、Refractive-Aware Gaussian Splatting (RA-GS)フレームワークを提案した。
水中の真の位置にある3D Gaussianを空気中での視点からの見かけの位置へとマッピングする微分可能な変換を導出することで、経験的な深度補正やデータ駆動型の事前分布に依存することなく、GSの効率的なラスタライザによるレンダリングから屈折の影響を受けない3次元再構成を実現する。
CGシミュレーションと、実環境でのフィールドワークで得られた河川環境データセットを用いた実験により、提案モデルが屈折歪みを除外し、高精度な画像合成と正確な幾何復元を達成することを示した。
シミュレーション環境においては定量的に、再構成された外観画像は屈折のない理想的な真値と比較して優れたSSIMなどの視覚品質を達成し、抽出された幾何形状はF1スコアが96\%を超え、Chamfer Distanceが0.011~mに達し、高い計測信頼性を確認した。
また、より複雑な事象を含む実環境でも、提案モデルが屈折の影響を除外し、mm級の再構成精度を達成することが示唆された。
これにより、空気中からの水中を対象とした二媒質水深写真測量において、屈折を正確なモデル化により除外するという実環境でも有効な手法を提案し、空白地帯となりやすかった浅水域において効率的かつ高精度な水深測量への実現へ大きく寄与すると考えている。

\section{Limitation and Future Work}\label{sec:limitation}

本研究における提案手法は、屈折を物理的な正確なモデル化で表現することで、高い性能を発揮するが、
前提として平坦な水面を考えており、波が立っている状況、すなわち水面の法線の方向が一定でない状況は考慮していない。
飛行高度や目標の測量精度に対応する空間解像度と、撮像において波の与える影響の大きさ(波の空間的周期や水深)が小さい場合、それらは平均的な誤差として無視できるかもしれないが、それには波と再構成の結果を対応つけるシミュレーションにおける実験が必要である。
また、例えカメラアライメントが正確であっても、波によってピクセル単位でシーンの一致しなくなるため、ピクセル単位でのシーンの一致を目的とするPSNRではなく、視覚的一致性を重要視するLIPISやSSIMをRA-GSの最適化におけるLoss関数に置くことで、影響を低減することができるかもしれない。
また、2DGSの法線一致やDepth Distortionなどの表面再構成を推進する正則化項\cite{Huang2024SIGGRAPH_2DGS,Gu2024CVPR_SuGaR,fan2024_trimGS}も、波の影響を低減させるのに有効であると考えられ、これらの影響の検証は今後の課題である。

\note{表: 平坦な水面の過程と= 波の不考慮}

また、本研究では、反射と屈折の割合を支配するフレネルの法則を含む、光の水面反射をモデル化していない。
フレネルの法則による空中角に増大に伴う、反射光の減衰は明示的な式で用意に微分可能なフレームワークとして導入可能であるが、その際に増加する外界の反射光の表現が難題になる。
本研究では、シミュレーション環境では反射の影響を無視し、実環境では偏光フィルターで反射の影響の低減を試みたが、その物理的な取り扱いは考慮していない。
特に小規模環境に対すて天空の水面反射は常に入射方向が一定となるため、球面調和関数による視線依存では表現しきれず、直接的な環境マップとしての外界の放射輝度を推定する必要がある。
幸いに、物体表面での反射に焦点を当て、環境光をGaussian Splattingのシーンと同時に推定する研究は複数提案されており、
\note{Geometry-Aware Reflection Disentanglement、3D Gaussian Splatting with Deferred Reflection、TransparentGS、SpecTRe-GS、Ref-GS : Directional Factorization for 2D Gaussian Splatting、GaussianShaderなど}がある。
例えば、GaussianShaderでは、フォトリアルなライティングをシミュレートするための微分可能な環境マップが提案されており、 
これらの手法を応用することで、環境光を同時に推定することが可能である。
\note{Cursorで実装を抽出してもらう。もう3年前の研究かい}

これらの効果は実環境の画像に本質的に含まれるため、明示的に扱わない限り、放射測光最適化においてノイズとして作用する。

水中での光の減衰も大きな問題となる。
これは、水媒質中で光の粒子が他の粒子と衝突する散乱、また吸収されることで生じる波長依存の効果であり、水中が青みがかって観測される物理的な理由である。
本研究では、実環境中では、水深が浅く無視できたが、水深が深くなるにつれて光の減衰は大きくなり、最終的に水底の物体から得られる放射輝度はなくなり、水中の散乱光と水面での反射光のみが観測されるようになる。
これらの物理的明示的なモデル化により、水中写真測量の水底が写っている範囲のみ再構成が可能なるという条件が、結果として測量の実行結果に反映され、水底が正確に視認できるという曖昧性を持った基準でなく、最大限の水深測量箇所が明確に反映されるようになる。
光の水中における減衰は、水中探索への重要性\cite{Trethewey2025_book-bathymetry}やそれを可能にする遠隔海洋探索マシンの導入などから、コンピュータビジョンでも重要な課題となっており盛んに研究されている。
特にSeaThru\cite{Akkaynak2019CVPR_Sea-Thru}では、それまでの霧などの大気中での一律な減衰モデルを改善し、RGBによる波長依存性を考慮した水中減衰モデルとその推定を定式化した。
WaterSplatting\cite{li20243DV_watersplatting}やSeaSplat\cite{Yang2025ICRA_SeaSplat}では、媒質中でその減衰を明示的にモデル化し、水質に依存する減衰係数の推定と媒質フリーな三次元再構成を実現した。
これらの研究は、本研究にも応用可能であり、水の濁りを除去し復元された色情報を含む水底の詳細な視覚情報は、エコトーンにおける生物の生息状況把握などに有用なアプリケーションを提供すると考えれる。

RA-GSでは、他の他視点三次元再構成手法と同様に、カメラパラメータを入力に必要とする。
本研究では実環境における解析では、水面をマスクし、陸域のStaticなSceneのみを用いて特徴点検出から始まるSfMプロセスを行うことで、正確なカメラパラメータと、歪みのない画像を得た。
実環境において大規模に、高い解像度で、陸域を十分に撮像できず、後続の水面マスク抽出手法が適用できない場合は、R-SfM\cite{Makris2024_refractive-aware-sfm}を用いて、屈折領域込みでのカメラポーズ推定が可能である。
また、RA-GSでは、陰関数定理による位置補正\cref{sec:position-correction}の実装によってカメラ外部パラメータに関しての勾配も追跡可能で、Photometricな情報からのカメラポーズ推定も可能であり、予め使用するカメラの内部キャリブレーションを行っておくことで、ジオリファレンス情報や屈折により誤差の増加したSfMによるカメラ位置姿勢を初期値として、シーンと同時にカメラ姿勢を最適化することも理論的に可能であり、実装と検証は今後の課題である。

本研究では、水面からのカメラ高度と水の屈折率を既知として、前者はSfMの再構成結果から、手動で水面がxy平面に平行になるように座標変換し、水面の位置を取得点群から推測し、後者は一般的な水の屈折率として一般的な1.33を用いた。
前述のように、RA-GSによる勾配追跡により水面からのカメラ高度は、カメラ外部座標の最適化問題に含有して推定可能である。
水の屈折率も水質により多少の変動があるが、これも最適化対象として理論的に推定可能であるが、最適化対象が多くなりすぎることで、安定した収束が可能になるか、実装に工夫と、詳細な検証、学習率設定が必要である。

本手研究ではオフナディア角を設けシーンを観測するという条件で、検証を行った。
この際、多視点性を確保する、一方で屈折により歪みの影響は増大するため(\cref{fig:space_compression_by_view})、屈折の影響を物理的に考慮できるRA-GSの有効性を発揮できる。
一方、実用においては、大規模かつ効率的なUAV写真測量のためには、特定の地点を中心としたオフナディア角を全域で行うのはデータ収集、解析における計算資源の観点から効率性を欠くため、オーバーラップ率を適当に設定した効率的なデータ収集が求められる。
この際には、一点を撮像する画像枚数と、その多視点性は減少するため、新規視点合成における少ない画像からの三次元再構成(Few-Shot Reconstruction)に近い問題設定になる。
この際、\~などのFew-Shot Gaussian Splatting手法や、前述の表面再構成を推進する正則化項の活用が重要になると考えられる。
\note{
  Depth-Regularized Optimization for 3D Gaussian Splatting in Few-Shot Images、
  FSGS: Real-Time Few-shot View Synthesis using Gaussian Splatting、
  FewViewGS: Gaussian Splatting with Few View Matching and Multi-stage Training
  }

大規模
広範囲のシーンを再構成するには、多くの画像と、シーンを表現するための多量のGaussianを配置する必要があり、計算コスト、メモリ使用量、データ容量の問題が生じる。
本研究では、宇治川Area1のデータセットにおいて、1/4にダウンサンプル後の概ね1300 $\times$ 1000 pixelsの画像、200枚を用いて、球面調和次数3で3M個のGaussianの最適化がメモリ上の限界であった(VRAM 12GB搭載のNVIDIA RTX 3090)。
これは、大規模シーンを効率的に復元する上で大きな制約となる。
こうした問題は、Gaussian Splattingの研究で実世界でのアプリケーションを実現する上で重要であり、盛んに研究されている。
大規模シーンの復元やメモリ効率を向上させるための、
\note{
  Horizon-GS: Unified 3D Gaussian Splatting for Large-Scale Aerial-to-Ground Scenes,
  A Hierarchical 3D Gaussian Representation for Real-Time Rendering of Very Large Datasets
}
や、
大きなデータ容量を必要とする再構成後のGaussianモデルを圧縮する
\note{
  Compact 3D Gaussian Representation for Radiance Field,
  3DGS. zip: A survey on 3D Gaussian Splatting Compression Methods
}
がある。

また、本研究では、実環境におけるワークフローでGaussianの初期化にMVSを用いた手法を提案したが、MVSには多くの計算時間と計算資源を必要となる。
そのため、初期化に依存しないMCMC Gaussian Splattingや、
Feed Forwardモデルからおおよその初期化を行う方法が大体として考えられる。

以上のように、Gaussian Splattingを始めとする、逆レンダリングを活用した三次元再構成、また、Feed Forwardな手法は現在も日進月歩の進歩を遂げており、それらの手法をRA-GSと組み合わせることで、水深写真測量の完全な実用化を推進できると考えている。

実用化においては、カメラポーズを正確に取得する必要があり、Structure from Motion (SfM) を用いることが標準的である。
本研究ではシミュレーションデータを用い、正確なカメラ内部・外部パラメータが既知であることを仮定して検証を進めたが、水面を含むシーンでは光の直進性が崩壊するため、実環境で正確なSfMの結果を得ることは困難である。
SfM内で屈折を考慮する研究~\cite{Makris2024_refractive-aware-sfm}では、カメラポーズが既知と仮定されているが、本研究で導入した座標変換は微分可能な変換に過ぎず、Bundle Adjustmentの過程に容易に組み込むことができる。
したがって、提案手法を応用することで、水面を写した画像から正確なカメラポーズ推定が可能になる可能性がある。
また、本研究ではカメラ高度から水面までの高さも既知と仮定していたが、実際の測量ではUAVの飛行高度設定によっておおよその推定・制御が可能である。
\note{キャリブレーションの制限は、ワークフローで前述したが、ここで述べる方がまとまりがある。}

水面からの相対的な高さをバンドル調整に組み込むことで、正確な水面位置の推定も可能になると考えている。

総括すると、本研究は写真測量バソメトリーにおける屈折歪みという長年の課題を解決するための重要な一歩を踏み出した。
光学的物理と明示的な3次元表現を密接に結合させることで、提案手法は浅水域における正確な深度マッピングの基盤を築き、屈折媒質における物理的基盤に基づいた再構成の新たな方向性を開拓した。
